\section{Setting, notation and the exact law}\label{sec:model}

\subsection{Setting and notation}\label{sec:notation}

\emph{The chain.} Fix $d\ge2$ and a generator $Q$ on $[d]=\{1,\dots,d\}$. Its
rates are $q_{ij}\ge0$ for $i\ne j$, the exit rates are $q_i=\sum_{j\ne i}q_{ij}$,
$Q_{ii}=-q_i$, and $q_{\max}=\max_iq_i$. When we write $q$ for the vector
$(q_i)$ this is clear from the context. (The scalar $q=1-p$ of Kagey's board
appears only where we recall Paper~A, and the scalar $q$ of
Proposition~\ref{prop:pair} is a common exit rate.) The start
$\mu$ is a probability law on $[d]$, and $\mu(F)=\sum_{i\in F}\mu_i$. For $N\ge2$
and $T>0$ put
\[
 h=\frac TN,\qquad L=\log N,\qquad \tau=\frac TL .
\]
The chain $X_1,\dots,X_N$ has $X_1\sim\mu$ and transition matrix $I+hQ$. So at
each step it jumps from $i$ to $j$ with probability $hq_{ij}$, and $N$ steps
look like the continuous-time chain with generator $Q$ run for time $T$. We
assume $hq_{\max}\le\frac12$ unless we say otherwise. The \emph{composition} is
the vector of visit counts $n=(n_1,\dots,n_d)$, with
$n_i=\#\{t\le N:X_t=i\}$, so that $n_1+\dots+n_d=N$. We write $P_N(n)$ for the
probability that the chain has composition $n$, $\supp n=\{i:n_i>0\}$ for its
support, and $\alpha=n/N$ for the visit fractions.

\emph{Faces.} A face is a nonempty set $F\subseteq[d]$, of size
$k=k_F=\lvert F\rvert$. Let $\Delta_F$ be the set of probability vectors on
$[d]$ that vanish off $F$, and $\Delta_F^\circ$ the set of those that are
positive on $F$. We number the states of $F$ as $1,\dots,k$ and use the
coordinates $(\alpha_2,\dots,\alpha_k)$ on $\Delta_F$ and $(t_2,\dots,t_k)$ on
$\{t\in\R^F:t_1+\dots+t_k=T\}$. We write $x\cdot y=\langle x,y\rangle
=\sum_{i\in F}x_iy_i$ and $x\circ y=(x_iy_i)_{i\in F}$. Let
$A_F=(q_{ij})_{i,j\in F}$ with zero diagonal, and $Q_F=A_F-\operatorname{diag}
(q_i)_{i\in F}$. Since $q_i$ includes the rates of leaving $F$, $Q_F$ generates
the chain killed when it leaves $F$. The switching graph $G_F$ has an edge
$i\to j$ whenever $i,j\in F$, $i\ne j$ and $q_{ij}>0$.

\emph{Perron roots.} A square matrix $B$ is Metzler if its off-diagonal entries
are nonnegative, and then its Perron root $\varrho(B)$ is its largest real
eigenvalue.
If $B$ is irreducible, $\varrho(B)$ is simple and has positive left and right
eigenvectors~\cite{bermanplemmons1994}.

\begin{definition}[faces]\label{def:faces}
Let $F$ be a face. We call $F$ \emph{irreducible} if $G_F$ is strongly connected,
so every singleton is irreducible. We call $F$ \emph{admissible} if for some
$N$ some path $x_1\cdots x_N$ of positive probability has support exactly $F$
(for $0<hq_{\max}<1$ this does not depend on $h$).
We say that $F$ satisfies \textup{(H)} if $F$ is irreducible and $\mu(F)>0$.
The \emph{exit rate} of $F$ is $\lambda_F=-\varrho(Q_F)$, so
$\lambda_{\{i\}}=q_i$. If $F$ is irreducible, $r_F$ and $l_F$ are right and left
eigenvectors of $Q_F$ for $-\lambda_F$ with positive entries and
$\langle l_F,r_F\rangle=1$, and $\alpha^*_F=l_F\circ r_F$ is the
\emph{quasi-ergodic law} of $F$~\cite{darrochseneta1965,darrochseneta1967}. The
\emph{face maximum} is $M_F=\max\{P_N(n):\supp n=F\}$, with $\max\emptyset=0$.
\end{definition}

We call $\lambda_F$ the exit rate for the following reason. If $F$ is
irreducible, then $e^{TQ_F}r_F=e^{-T\lambda_F}r_F$ and $r_F$ is positive. So
the chance that the continuous-time chain started at a state of $F$ stays in
$F$ up to time $T$, which is a row sum of $e^{TQ_F}$, lies between
$e^{-T\lambda_F}\min r_F/\max r_F$ and $e^{-T\lambda_F}\max r_F/\min r_F$. If
$F$ is reducible, we order the strongly connected components $C_1,\dots,C_r$ of
$G_F$ so that every edge between components goes forward. Then $Q_F$ is block
triangular with diagonal blocks $Q_{C_j}$, and
$\lambda_F=\min_j\lambda_{C_j}$. In all cases $0\le\lambda_F\le\min_{i\in
F}q_i$, since the rows of $Q_F$ have sums at most $0$ and a Perron root is at
least every diagonal entry. For irreducible $F$ the quasi-ergodic law
$\alpha^*_F$ is the long-run law of the visit fractions of the chain conditioned
to stay in $F$, and the best composition on $F$ sits near $N\alpha^*_F$
(Theorem~\ref{thm:facemax}).

Later sections add their own symbols where they are needed, namely the rate
$I_F$ (Section~\ref{sec:firstorder}), the tilted chain and the constant $K_F$
(Definition~\ref{def:tilt}), and the planar notation of
Section~\ref{sec:planar}.

\emph{Kagey's board.} For $a,b\ge1$ let $S_{a,b}(u)=\sum_wu^{s(w)}$, the sum
over the words $w$ with $a$ right steps and $b$ left steps, where $s(w)$ is the
number of switches of $w$. Each such word has probability
$\frac12p^{N-1}(q/p)^{s(w)}$ on Kagey's board with $N=a+b$ steps. So
$S_{a,b}(q/p)$ is the probability of bin $a$ divided by the probability
$\frac12p^{N-1}$ of an endpoint. For $N\ge2$ let
$S_N=S_{\lfloor N/2\rfloor,\lceil N/2\rceil}$, which belongs to a middle bin.
It is a nonzero polynomial with nonnegative coefficients
and no constant term, so $S_N(u)=1$ has exactly one positive root $u_N$, and
$p_N=1/(1+u_N)$ is the crossing of the middle and the endpoints. Put
$z_N=N(1-p_N)$ and $c_0=\frac12\log(\pi/8)$. Paper~A proves the crossing
equation $z_N+\frac12\log z_N=L+c_0+1/(8z_N)+R_N$ with an explicit remainder
$R_N=O(1/L^2)$.

\emph{Conventions.} Unless we say otherwise, $Q$, $\mu$ and the faces under
discussion are fixed, and every constant, every implied constant in $O(\cdot)$
and every threshold $N_0$ depends only on them and on the compact set
$\mathcal K$ when one appears. A bound that holds ``uniformly for $1\le T\le
N^{1/3}$'' holds for all $N\ge N_0$ and all $T$ in that range with one implied
constant. Subscripts as in $O_{\mathcal K}$ stress a dependence.

\subsection{Admissible faces}

\begin{lemma}[admissible faces]\label{lem:admissible}
Let $F$ be a face and $C_1,\dots,C_r$ the strongly connected components of
$G_F$. \textup{(a)} $F$ is admissible if and only if the components can be
numbered so that $\mu(C_1)>0$ and, for each $j<r$, $q_{xy}>0$ for some $x\in
C_j$ and $y\in C_{j+1}$. \textup{(b)} An irreducible face is admissible if and
only if $\mu(F)>0$. So \textup{(H)} holds exactly for the irreducible admissible
faces. \textup{(c)} If $Q$ has symmetric support, that is $q_{ij}>0$ exactly
when $q_{ji}>0$, then every admissible face is irreducible.
\end{lemma}

\begin{proof}
We translate paths into walks in $G_F$. Take a path of positive probability
with support $F$ and list the states of its maximal runs. This list is a walk
in $G_F$ that starts at a state of $\supp\mu$ and visits every state of $F$.
Conversely, every such walk is the list of runs of a path of positive
probability with support $F$, as soon as $N$ is at least the length of the
walk. (The path follows the walk and then stays in its last state, and a stay
has positive probability.) So $F$ is admissible exactly when $G_F$ has a walk
that starts in $\supp\mu$ and visits every state of $F$.

(a) Suppose that there is such a walk. Once it has left a component $C$ for a
state $y\notin C$, it cannot come back to $C$. (If it did, $y$ could be reached
from $C$ and $C$ from $y$, so $y$ would belong to $C$.) Hence the walk visits
each component in one block of consecutive times, and it passes from one block
to the next along a single edge. We number the components in the order of these
blocks. The walk starts in $C_1$, so $\mu(C_1)>0$, and the edge from block $j$
to block $j+1$ is an edge $x\to y$ with $x\in C_j$ and $y\in C_{j+1}$. This is
the condition in (a).

Conversely, suppose that the components are numbered as in (a), with edges
$x_j\to y_{j+1}$ from $C_j$ to $C_{j+1}$ for $j<r$. Take $y_1\in C_1$ with
$\mu_{y_1}>0$ and any $x_r\in C_r$. Since $C_j$ is strongly connected, there is
a walk inside $C_j$ from $y_j$ to $x_j$ that visits every state of $C_j$. We
join these $r$ walks by the edges $x_j\to y_{j+1}$. The result is a walk that
starts in $\supp\mu$ and visits every state of $F$.

(b) This is (a) with $r=1$.

(c) If $Q$ has symmetric support, an edge from one component to another would
come with an edge back, and the two components would be the same. So no edge
joins two components, and the condition in (a) forces $r=1$.
\end{proof}

So $M_F>0$ for large $N$ if $F$ is admissible, and $M_F=0$ otherwise.

\subsection{The exact law}

For $m\in\Z_{\ge1}^F$ put $M=\sum_{i\in F}m_i$ and
$W(m)=\sum_c\mu_{c_1}\prod_{j<M}q_{c_jc_{j+1}}$, the sum over the words
$c=c_1\cdots c_M$ in $F$ with $c_j\ne c_{j+1}$ in which each $i\in F$ occurs
$m_i$ times. Such a word lists the states of the runs of a path, and $m_i$ is
the number of runs in state $i$.

\begin{lemma}[run expansion]\label{lem:runs}
Let $hq_{\max}\le1$ and let $n$ have support $F$. Then, with
$\binom{n_i-1}{m_i-1}=0$ for $m_i>n_i$ and $0^0=1$,
\[
P_N(n)=\sum_{m\in\Z_{\ge1}^F}W(m)\,h^{M-1}\prod_{i\in F}\binom{n_i-1}{m_i-1}
(1-hq_i)^{n_i-m_i}.
\]
\end{lemma}

\begin{proof}
We group the paths with composition $n$ by the word of their runs. The states
$c_1\cdots c_M$ of the maximal runs of a path form a word as in $W(m)$, where
$m_i$ is the number of runs in state $i$. The path starts at $c_1$, which has
probability $\mu_{c_1}$. It has $M-1$ switches, of probabilities
$hq_{c_jc_{j+1}}$, and $n_i-m_i$ stays in state $i$, of probability $1-hq_i$
each. So every path with the word $c$ and composition $n$ has the same
probability, namely
$\mu_{c_1}\prod_{j<M}q_{c_jc_{j+1}}\,h^{M-1}\prod_i(1-hq_i)^{n_i-m_i}$. The
number of these paths is $\prod_i\binom{n_i-1}{m_i-1}$, because the lengths of
the $m_i$ runs in state $i$ can be any of the $\binom{n_i-1}{m_i-1}$ ways to
write $n_i$ as an ordered sum of $m_i$ positive parts. Summing over the words
$c$ with given $m$ gives $W(m)$, and summing over $m$ gives the claim.
\end{proof}

Lemma~\ref{lem:runs} is checked in Lean (Section~\ref{sec:verify}). For
$F=\{i\}$ the only path is $i^N$, so uniformly for $1\le T\le N^{1/2}$
\begin{equation}\label{eq:vertexmax}
M_{\{i\}}=\mu_i(1-hq_i)^{N-1}=\mu_ie^{-q_iT}\bigl(1+O(T^2/N)\bigr).
\end{equation}
For $t\in(0,\infty)^F$ let $H(t)=\sum_mW(m)\prod_{i\in F}t_i^{m_i-1}/(m_i-1)!$,
which is finite since $W(m)\le k^Mq_{\max}^{M-1}$. The next lemma says that
$e^{-\langle q,t\rangle}H(t)$ is the density of the occupation times of the
continuous-time chain on the event that it stays in $F$. It is the continuum
form of Lemma~\ref{lem:runs}, in which the binomial coefficients become the
powers $t_i^{m_i-1}/(m_i-1)!$. Lemma~\ref{lem:discrete} below shows that the
discrete law is close to it.

\begin{lemma}[the continuum density]\label{lem:continuum}
Let $(Y_s)_{s\ge0}$ be the continuous-time chain with generator $Q$ and
$Y_0\sim\mu$, let $\mathcal O_i=\int_0^T\mathbf 1\{Y_s=i\}\,ds$, and let $E_F$
be the event that $Y$ stays in $F$ during $[0,T]$ and visits every state of $F$.
Then $\Prob(E_F,\ \mathcal O\in\cdot\,)$ has the density
$e^{-\langle q,t\rangle}H(t)$ in $(t_2,\dots,t_k)$ on
$\{t\in(0,\infty)^F:\sum_it_i=T\}$, and $\Prob(E_F,\ \mathcal O/T\in\cdot\,)$
has the density
\[
f_T(\alpha)=T^{k-1}e^{-T\langle q,\alpha\rangle}H(T\alpha)
\]
in $(\alpha_2,\dots,\alpha_k)$. For $n$ with support $F$ and $t=hn$,
$N^{-(k-1)}f_T(n/N)=h^{k-1}e^{-\langle q,t\rangle}H(t)$. For $F=\{i\}$,
$f_T=\Prob(E_F)=\mu_ie^{-q_iT}$.
\end{lemma}

\begin{proof}
We split $E_F$ by the sequence of states that $Y$ visits, and for each
sequence we integrate out the holding times. Up to a null set, $E_F$ is the
disjoint union of the events $E_c$ that $Y$ visits exactly $c_1,\dots,c_M$ in
this order during $[0,T]$, over the words $c$ in the sums $W(m)$.

\emph{The density of the holding times.} Fix $c$, and let $s_j$ be the time
that $Y$ spends in its $j$-th run during $[0,T]$, so that $s_1+\dots+s_M=T$.
For $j<M$ the $j$-th run is complete. Its holding time has the density
$q_{c_j}e^{-q_{c_j}s_j}$, and it ends with a jump to $c_{j+1}$, which has
probability $q_{c_jc_{j+1}}/q_{c_j}$. The last run only has to last longer
than $s_M=T-\sum_{j<M}s_j$, which has probability $e^{-q_{c_M}s_M}$. We
multiply these factors. On the set where $s_j>0$ for every $j\le M$, the event
$E_c$ has the density
\[
\mu_{c_1}\prod_{j<M}q_{c_jc_{j+1}}\;e^{-\langle q,\mathcal O\rangle}
\]
in $(s_1,\dots,s_{M-1})$, where $\mathcal O_i$ is the sum of the $s_j$ with
$c_j=i$.

\emph{From holding times to occupation times.} This density depends on the
holding times only through $\mathcal O$, so we change variables. For each state
$i\ne c_M$ we replace the last holding time in state $i$ by $\mathcal O_i$.
This change of variables is triangular with unit diagonal, so it preserves
Lebesgue measure. The new variables are the $k-1$ occupation times
$\mathcal O_i$ with $i\ne c_M$, and for each state $i$ the other $m_i-1$
holding times in state $i$ (for $i=c_M$, those before the last run). Now fix
$\mathcal O=t$. The $m_i-1$ free holding times in state $i$ are positive and
their sum is less than $t_i$, so they range over a simplex of volume
$t_i^{m_i-1}/(m_i-1)!$. (For two states and the word $c=121$ the only free
variable is the first holding time, which ranges over $(0,t_1)$, and the factor
is $t_1$.) Integrating the free holding times out, we find that
$\Prob(E_c,\ \mathcal O\in\cdot\,)$ has the density
$\mu_{c_1}\prod_{j<M}q_{c_jc_{j+1}}\,e^{-\langle q,t\rangle}
\prod_{i\in F}t_i^{m_i-1}/(m_i-1)!$ in the coordinates $t_i$ with $i\ne c_M$.

\emph{Summing over the words.} On $\{\sum_it_i=T\}$ any $k-1$ of the
coordinates $t_i$ give the same Lebesgue measure, so we may use
$(t_2,\dots,t_k)$ for every $c$. Summing over the words $c$ with given $m$
gives $W(m)$, and summing over $m$ gives $e^{-\langle q,t\rangle}H(t)$. The
substitution $t=T\alpha$ has Jacobian $T^{k-1}$, which gives $f_T$, and
$T/N=h$ gives the identity for $t=hn$. For $F=\{i\}$ the event $E_F$ says that
$Y$ starts at $i$ and does not jump before time $T$.
\end{proof}

\begin{example}[the four models]\label{ex:models}
(i) For $d=2$, $q_{12}=q_{21}=1$ and uniform $\mu$ the chain is Kagey's board,
the persistent walk of~\cite{paperA}, with $p=1-h$. So $T=N(1-p)$, the faces
$\{1\}$ and $\{2\}$ are the endpoints, and Lemma~\ref{lem:runs} is the run count
of~\cite{paperA}. (ii) For $Q=J-dI$, with $J$ the all-ones matrix, and uniform
$\mu$ it is the model of~\cite{paperB}, which switches to each other state with
probability $h$. (iii) For $Q=\Pi-I$, with $\Pi$ stochastic with zero diagonal,
$\lambda_F=1-\varrho(\Pi_F)$. This is the sticky chain of
Section~\ref{sec:sticky}. (iv) Let $d=4$, number the directions east, north,
west and south as $0,1,2,3$ modulo $4$, and put $q_{j,j\pm1}=r$, $q_{j,j+2}=s$
and $\mu$ uniform. Then $X_t$ is the direction of step $t$ of a planar walk that
turns left or right with probability $rh$ each and reverses with probability
$sh$ (Section~\ref{sec:planar}). For $s=0$ and $r=1$ this is the 4-cycle of
Section~\ref{sec:worked}. In (i), (ii) and (iv), $Q$ has symmetric support and
$\mu$ full support, so all admissible faces satisfy \textup{(H)} by
Lemma~\ref{lem:admissible}.
\end{example}
